% ============================================================
%  CHAPTER 4 – INTERFACES
% ============================================================
\chapter{Interface Control}
\label{ch:interface-control}
This chapter describes each of the interfaces introduced in chapter~\ref{ch:system-design}. Each interface listed in section~\ref{sec:summary-of-interfaces}, is described and allocated to a role in section~\ref{sec:work-allocation}.

\section{Work Allocation}
\label{sec:work-allocation}
This section ascribes each interface sa the responsibility of the role under which it is described. Each member read and confirmed their responsibilities and interfaces under their jurisdiction as indicated below. By having their name printed next to their role they have accepted the responsibility and are the liable person for all interfaces listed in their section.

\subsection{EERI I:  }
\label{subsec:eeri1-interfaces}
EERI I Drive Electroncics by Jopie Roux.
\begin{table}[H]
    \centering
    \caption{Operational Interfaces.}
    \label{tab1:interfaces-eeri1}
    \begin{tabular}{@{} l l l p{5cm} @{}}
        \toprule
        \textbf{Interface} & \textbf{From} & \textbf{To} & \textbf{Description} \\
        \midrule
        I1 & F/U 1.0 & F/U 2.0 & Description of data/power/signal. \\
        I/F 1.2 & F/U 2.0 & F/U 3.0 & Description. \\
        I/F 1.2 & F/U 2.0 & F/U 3.0 & Description. \\
        I/F 1.2 & F/U 2.0 & F/U 3.0 & Description. \\
        I/F 1.2 & F/U 2.0 & F/U 3.0 & Description. \\
        I/F 1.2 & F/U 2.0 & F/U 3.0 & Description. \\
        I/F 1.2 & F/U 2.0 & F/U 3.0 & Description. \\
        \bottomrule
    \end{tabular}
\end{table}


\begin{table}[H]
    \centering
    \caption{Operational Interfaces.}
    \label{tab2:interfaces-eeri1}
    \begin{tabular}{@{} l l l p{5cm} @{}}
        \toprule
        \textbf{Interface} & \textbf{From} & \textbf{To} & \textbf{Description} \\
        \midrule
        I1 & F/U 1.0 & F/U 2.0 & Description of data/power/signal. \\
        I/F 1.2 & F/U 2.0 & F/U 3.0 & Description. \\
        I/F 1.2 & F/U 2.0 & F/U 3.0 & Description. \\
        I/F 1.2 & F/U 2.0 & F/U 3.0 & Description. \\
        I/F 1.2 & F/U 2.0 & F/U 3.0 & Description. \\
        I/F 1.2 & F/U 2.0 & F/U 3.0 & Description. \\
        I/F 1.2 & F/U 2.0 & F/U 3.0 & Description. \\
        \bottomrule
    \end{tabular}
\end{table}

\subsection{EERI II: }
\label{subsec:eeri2-interfaces}
EERI II Magnetic Load Attachment by JP Smith.
\begin{table}[H]
    \centering
    \caption{Operational Interfaces.}
    \label{tab:interfaces-eeri2}
    \begin{tabular}{@{} l l l p{5cm} @{}}
        \toprule
        \textbf{Interface} & \textbf{From} & \textbf{To} & \textbf{Description} \\
        \midrule
        I1 & F/U 1.0 & F/U 2.0 & [Description of data/power/signal.] \\
        I2 & F/U 2.0 & F/U 3.0 & [Description.] \\
        \bottomrule
    \end{tabular}
\end{table}



\subsection{MCTR I: }
\label{subsec:mctr1-interfaces}
MCTR I Lateral Movement by Romeo Mapasa.
\begin{table}[H]
    \centering
    \caption{Operational Interfaces.}
    \label{tab:interfaces-mctr1}
    \begin{tabular}{@{} l l l p{5cm} @{}}
        \toprule
        \textbf{Interface} & \textbf{From} & \textbf{To} & \textbf{Description} \\
        \midrule
        I1 & F/U 1.0 & F/U 2.0 & [Description of data/power/signal.] \\
        I2 & F/U 2.0 & F/U 3.0 & [Description.] \\
        \bottomrule
    \end{tabular}
\end{table}




\subsection{MCTR II: Juan Fourie}
\label{subsec:mctr2-interfaces}
MCTR II Spooling by Juan Fourie.
% =====================================================================
%  mctr2_parameters.tex  --  MCTR II (Spooling) parameter + calculation sheet
%
%  HOW TO USE
%   * Edit ONLY the \In... values in Section 1. Everything in Section 2
%     (\Raw... results) is computed from them by \fpeval, so every
%     table, equation and sentence that quotes a result updates on the
%     next compile.
%   * Units are noted per line. Torques in N.m unless stated.
%   * Display helpers:  \shown{Name}{digits}        -> number only
%                       \showv{Name}{digits}{unit}  -> number + unit
%   * This file is safe to \input from several chapters (loads once).
%   * Needs: LaTeX >= 2021 (\fpeval in kernel), amsmath, xcolor.
% =====================================================================
\ifdefined\MCTRIIparamsloaded\else
\def\MCTRIIparamsloaded{}

\providecommand{\tbd}[1]{\textcolor{red}{\textbf{[TBD: #1]}}}
\newcommand{\val}[2]{\ensuremath{#1\,\mathrm{#2}}}
\newcommand{\shown}[2]{\fpeval{round(\csname Raw#1\endcsname,#2)}}
\newcommand{\showv}[3]{\val{\fpeval{round(\csname Raw#1\endcsname,#2)}}{#3}}
% PASS/FAIL helpers: \passge{value}{limit} passes when value >= limit
\newcommand{\passge}[2]{\ifdim #1pt<#2pt \textcolor{red}{\textbf{FAIL}}\else\textcolor{green!45!black}{\textbf{PASS}}\fi}
\newcommand{\passle}[2]{\ifdim #1pt>#2pt \textcolor{red}{\textbf{FAIL}}\else\textcolor{green!45!black}{\textbf{PASS}}\fi}

% Functional-unit number of MCTR II (Level 1 drawing: F/U 1.4).
% NOTE: the Level 2 drawing currently labels sub-units 1.3.x; change the
% drawing to 1.4.x (or set this to 1.3 to match the drawing, but 1.3 is MCTR I).
\newcommand{\fuMII}{1.3}

% ---------------------------------------------------------------------
% 1. INPUTS  (edit these)
% ---------------------------------------------------------------------
% -- Loads, stroke, timing
\newcommand{\InMassRated}{20}    % kg   rated payload (PHY1)
\newcommand{\InMassDesign}{30}   % kg   design overload case (analysis only)
\newcommand{\InMassAttach}{2}    % kg   magnet + hook hanging on rope  [TBD: weigh]
\newcommand{\InGrav}{9.81}       % m/s^2
\newcommand{\InStroke}{1.6}      % m     hoist stroke (PT2)
\newcommand{\InTimeReq}{20}      % s     max stroke time (PT2)
\newcommand{\InServiceF}{1.5}    % -     service factor on drum torque
\newcommand{\InSFreq}{1.5}       % -     minimum acceptable safety factor (strength checks)
% -- Drum and rope
\newcommand{\InDrumD}{50}        % mm    drum core diameter (FROZEN; 70 mm entry removed)
\newcommand{\InDrumH}{60}        % mm    drum usable length
\newcommand{\InRopeD}{5}         % mm    rope diameter
\newcommand{\InRopeCentre}{1}    % 1 = tension acts at rope centre-line (D+d); 0 = old spreadsheet (D only)
% -- Motor and VSD
\newcommand{\InMotorP}{180}      % W     Askari MS632-4
\newcommand{\InMotorN}{1360}     % rpm   at rated frequency
\newcommand{\InMotorF}{50}       % Hz    rated frequency
\newcommand{\InVsdF}{50}         % Hz    frequency VSD will actually deliver  [TBD: EEII I to confirm]
% -- Gearbox (stage A = motor side, stage B = drum side)
\newcommand{\InNpA}{16}   \newcommand{\InNgA}{102}   % pinion / gear teeth, stage A
\newcommand{\InNpB}{16}   \newcommand{\InNgB}{102}   % pinion / gear teeth, stage B
\newcommand{\InModA}{1.5} \newcommand{\InModB}{2}    % mm module
\newcommand{\InFwA}{10}   \newcommand{\InFwB}{15}    % mm face width
\newcommand{\InPhi}{25}          % deg   pressure angle (use 24 if Inventor forces it; both gears of a mesh must match)
\newcommand{\InYpA}{0.296}       % Lewis form factor, pinion  (Shigley Tbl 14-2, 20 deg value: conservative for 25 deg)
\newcommand{\InYgA}{0.447}       % Lewis form factor, gear
\newcommand{\InYpB}{0.296}
\newcommand{\InYgB}{0.447}
\newcommand{\InEta}{0.95}        % -     efficiency per gear stage
\newcommand{\InSigAllow}{20}     % MPa   allowable PETG root stress [GUESS: replace with printed-coupon result / SF]
\newcommand{\InClear}{0.25}      % mm    added centre-distance clearance (printed backlash)
\newcommand{\InBed}{220}         % mm    print bed edge length [TBD: confirm printer]
% -- Shaft, hub and fastener
\newcommand{\InShaftD}{10}       % mm    shaft diameter
\newcommand{\InSpan}{100}        % mm    bearing span (centre-to-centre)
\newcommand{\InKt}{2}            % -     bending stress-concentration factor, cross-hole
\newcommand{\InKts}{2}           % -     torsion stress-concentration factor, cross-hole
\newcommand{\InSy}{250}          % MPa   yield strength  [TBD: replace with certified grade, e.g. EN8 bright bar]
\newcommand{\InBoltD}{4}         % mm    cross-bolt diameter
\newcommand{\InHubWall}{8}       % mm    hub wall length bearing on bolt
\newcommand{\InHubAllow}{20}     % MPa   allowable PETG bearing stress [GUESS]
% -- Brake
\newcommand{\InBrakeSF}{2}       % -     brake torque safety factor
\newcommand{\InBrakeRated}{0}    % N.m   datasheet holding torque; 0 = not yet known  [TBD: datasheet]
% -- Load cell and HX711
\newcommand{\InCellCap}{100}     % kg    YZC-516 rated capacity
\newcommand{\InCellSens}{2}      % mV/V  rated output [TBD: datasheet]
\newcommand{\InExc}{5}           % V     excitation (HX711 supply)
\newcommand{\InGain}{128}        % -     HX711 channel A gain
\newcommand{\InSps}{80}          % SPS   HX711 output rate (RATE pin high)
\newcommand{\InFilt}{4}          % -     samples averaged by interlock filter
\newcommand{\InNoise}{500}       % counts peak-to-peak noise [GUESS: measure with motor running]
\newcommand{\InTrip}{1.25}       % -     overload trip as multiple of (rated + attachment)
% -- Cost (R, incl. VAT and shipping, from Cost_and_Budget.xlsx)
\newcommand{\InBudget}{5000}
\newcommand{\InCostBearings}{116.24}
\newcommand{\InCostCell}{429.00}     % [CHECK: sheet unit 493.35 + 135.70 shipping = 629.05]
\newcommand{\InCostHX}{179.40}       % [CHECK: unit price 43.70]
\newcommand{\InCostPSU}{390.17}
\newcommand{\InCostBrake}{907.68}    % excluded from the sheet's own subtotal

% ---------------------------------------------------------------------
% 2. DERIVED RESULTS  (do not edit)
% ---------------------------------------------------------------------
% -- Kinematics
\newcommand{\RawDeff}{\fpeval{\InDrumD+\InRopeCentre*\InRopeD}}
\newcommand{\RawVreq}{\fpeval{\InStroke*1000/\InTimeReq}}
\newcommand{\RawWraps}{\fpeval{\InStroke*1000/(pi*\RawDeff)}}
\newcommand{\RawWrapsFit}{\fpeval{floor(\InDrumH/\InRopeD)}}
\newcommand{\RawiA}{\fpeval{\InNgA/\InNpA}}
\newcommand{\RawiB}{\fpeval{\InNgB/\InNpB}}
\newcommand{\RawiT}{\fpeval{\RawiA*\RawiB}}
\newcommand{\RawNmot}{\fpeval{\InMotorN*\InVsdF/\InMotorF}}
\newcommand{\RawNdrum}{\fpeval{\RawNmot/\RawiT}}
\newcommand{\RawVline}{\fpeval{pi*\RawDeff*\RawNdrum/60}}
\newcommand{\RawTstroke}{\fpeval{\InStroke*1000/\RawVline}}
\newcommand{\RawSpeedMargin}{\fpeval{(\RawVline/\RawVreq-1)*100}}
\newcommand{\RawiReq}{\fpeval{pi*\RawDeff*\RawNmot/(60*\RawVreq)}}
\newcommand{\RawFmin}{\fpeval{\RawVreq*60*\RawiT/(pi*\RawDeff)*\InMotorF/\InMotorN}}
% -- Drum torque, motor sizing
\newcommand{\RawFropeR}{\fpeval{\InMassRated*\InGrav}}
\newcommand{\RawFropeD}{\fpeval{\InMassDesign*\InGrav}}
\newcommand{\RawTdrR}{\fpeval{\InMassRated*\InGrav*\RawDeff/2000}}
\newcommand{\RawTdrD}{\fpeval{\InMassDesign*\InGrav*\RawDeff/2000*\InServiceF}}
\newcommand{\RawTdrSheet}{\fpeval{\InMassDesign*\InGrav*\InDrumD/2000*\InServiceF}}
\newcommand{\RawEtaGB}{\fpeval{\InEta*\InEta}}
\newcommand{\RawTmReq}{\fpeval{\RawTdrR/(\RawiT*\RawEtaGB)}}
\newcommand{\RawTmRated}{\fpeval{\InMotorP/(2*pi*\InMotorN/60)}}
\newcommand{\RawTmMargin}{\fpeval{\RawTmRated/\RawTmReq}}
\newcommand{\RawPmech}{\fpeval{\InMassRated*\InGrav*\RawVline/1000/\RawEtaGB}}
\newcommand{\RawPavail}{\fpeval{\InMotorP*\InVsdF/\InMotorF}}
% -- Gear loads (design case) and Lewis stress
\newcommand{\RawTtwo}{\fpeval{\RawTdrD/(\RawiB*\InEta)}}
\newcommand{\RawTone}{\fpeval{\RawTtwo/(\RawiA*\InEta)}}
\newcommand{\RawFtB}{\fpeval{2*\RawTtwo*1000/(\InModB*\InNpB)}}
\newcommand{\RawFtA}{\fpeval{2*\RawTone*1000/(\InModA*\InNpA)}}
\newcommand{\RawFnB}{\fpeval{\RawFtB/cosd(\InPhi)}}
\newcommand{\RawFnA}{\fpeval{\RawFtA/cosd(\InPhi)}}
\newcommand{\RawFrB}{\fpeval{\RawFtB*tand(\InPhi)}}
\newcommand{\RawSigPA}{\fpeval{\RawFtA/(\InFwA*\InModA*\InYpA)}}
\newcommand{\RawSigGA}{\fpeval{\RawFtA/(\InFwA*\InModA*\InYgA)}}
\newcommand{\RawSigPB}{\fpeval{\RawFtB/(\InFwB*\InModB*\InYpB)}}
\newcommand{\RawSigGB}{\fpeval{\RawFtB/(\InFwB*\InModB*\InYgB)}}
\newcommand{\RawSFpA}{\fpeval{\InSigAllow/\RawSigPA}}
\newcommand{\RawSFpB}{\fpeval{\InSigAllow/\RawSigPB}}
% -- Gear geometry
\newcommand{\RawDpinA}{\fpeval{\InModA*\InNpA}}
\newcommand{\RawDgearA}{\fpeval{\InModA*\InNgA}}
\newcommand{\RawDpinB}{\fpeval{\InModB*\InNpB}}
\newcommand{\RawDgearB}{\fpeval{\InModB*\InNgB}}
\newcommand{\RawDoGearA}{\fpeval{\InModA*(\InNgA+2)}}
\newcommand{\RawDoGearB}{\fpeval{\InModB*(\InNgB+2)}}
\newcommand{\RawCdA}{\fpeval{\InModA*(\InNpA+\InNgA)/2+\InClear}}
\newcommand{\RawCdB}{\fpeval{\InModB*(\InNpB+\InNgB)/2+\InClear}}
\newcommand{\RawZmin}{\fpeval{2/(sind(\InPhi)^2)}}
\newcommand{\RawDbPinB}{\fpeval{\RawDpinB*cosd(\InPhi)}}
\newcommand{\RawBedUse}{\fpeval{max(\RawDoGearA,\RawDoGearB)/\InBed*100}}
% -- Shaft (macros with arguments: #1 = payload mass kg, #2 = torque factor)
\newcommand{\Tdr}[2]{\fpeval{#1*\InGrav*\RawDeff/2*#2}}                                   % N.mm
\newcommand{\FgB}[2]{\fpeval{2*\Tdr{#1}{#2}/(\InEta*\InModB*\InNgB*cosd(\InPhi))}}        % N
\newcommand{\Msh}[2]{\fpeval{(#1*\InGrav+\FgB{#1}{#2})*\InSpan/4}}                        % N.mm
\newcommand{\SigBsh}[2]{\fpeval{32*\Msh{#1}{#2}/(pi*\InShaftD^3)}}                        % MPa
\newcommand{\Taush}[2]{\fpeval{16*\Tdr{#1}{#2}/(pi*\InShaftD^3)}}                         % MPa
\newcommand{\VMsh}[2]{\fpeval{sqrt((\InKt*\SigBsh{#1}{#2})^2+3*(\InKts*\Taush{#1}{#2})^2)}} % MPa
\newcommand{\RawFgD}{\FgB{\InMassDesign}{\InServiceF}}
\newcommand{\RawMshD}{\Msh{\InMassDesign}{\InServiceF}}
\newcommand{\RawSigBD}{\SigBsh{\InMassDesign}{\InServiceF}}
\newcommand{\RawTauD}{\Taush{\InMassDesign}{\InServiceF}}
\newcommand{\RawVMD}{\VMsh{\InMassDesign}{\InServiceF}}
\newcommand{\RawVMR}{\VMsh{\InMassRated}{1}}
\newcommand{\RawSFsD}{\fpeval{\InSy/\RawVMD}}
\newcommand{\RawSFsR}{\fpeval{\InSy/\RawVMR}}
% -- Hub / cross-bolt (design case, drum shaft)
\newcommand{\RawFside}{\fpeval{\RawTdrD*1000/\InShaftD}}
\newcommand{\RawSigHub}{\fpeval{\RawFside/(\InBoltD*\InHubWall)}}
\newcommand{\RawTauBolt}{\fpeval{\RawFside/(pi*\InBoltD^2/4)}}
\newcommand{\RawSFhub}{\fpeval{\InHubAllow/\RawSigHub}}
% -- Rated-case (20 kg, no service factor) checks for gears and hub
\newcommand{\RawLoadRatio}{\fpeval{\RawTdrR/\RawTdrD}}
\newcommand{\RawSFpAR}{\fpeval{\RawSFpA/\RawLoadRatio}}
\newcommand{\RawSFpBR}{\fpeval{\RawSFpB/\RawLoadRatio}}
\newcommand{\RawSigHubR}{\fpeval{\RawSigHub*\RawLoadRatio}}
\newcommand{\RawSFhubR}{\fpeval{\InHubAllow/\RawSigHubR}}
% -- Brake
\newcommand{\RawTholdDrum}{\fpeval{\InMassDesign*\InGrav*\RawDeff/2000}}
\newcommand{\RawTholdIn}{\fpeval{\RawTholdDrum/\RawiT}}
\newcommand{\RawBrakeReqDrum}{\fpeval{\RawTholdDrum*\InBrakeSF}}
\newcommand{\RawBrakeReqIn}{\fpeval{\RawTholdIn*\InBrakeSF}}
\newcommand{\BrakeMarginIn}{\ifdim\InBrakeRated pt=0pt\tbd{datasheet torque}\else\fpeval{round(\InBrakeRated/\RawBrakeReqIn,1)}\fi}
\newcommand{\BrakeMarginDrum}{\ifdim\InBrakeRated pt=0pt\tbd{datasheet torque}\else\fpeval{round(\InBrakeRated/\RawBrakeReqDrum,1)}\fi}
% -- Load cell and HX711
\newcommand{\RawFcell}{\fpeval{(\InMassRated+\InMassAttach)*\InGrav}}
\newcommand{\RawMtrip}{\fpeval{(\InMassRated+\InMassAttach)*\InTrip}}
\newcommand{\RawCellUse}{\fpeval{\RawMtrip/\InCellCap*100}}
\newcommand{\RawSigFS}{\fpeval{\InCellSens*\InExc}}
\newcommand{\RawAdcRange}{\fpeval{0.5*\InExc/\InGain*1000}}
\newcommand{\RawCntPerKg}{\fpeval{\InCellSens/1000/\InCellCap*\InGain/0.5*2^23}}
\newcommand{\RawResG}{\fpeval{\InNoise/\RawCntPerKg*1000}}
\newcommand{\RawLatency}{\fpeval{\InFilt/\InSps*1000}}
\newcommand{\RawTravel}{\fpeval{\RawVline*\RawLatency/1000}}
% -- Cost
\newcommand{\RawCostTot}{\fpeval{\InCostBearings+\InCostCell+\InCostHX+\InCostPSU+\InCostBrake}}
\newcommand{\RawBudgetLeft}{\fpeval{\InBudget-\RawCostTot}}
\newcommand{\RawBudgetUse}{\fpeval{\RawCostTot/\InBudget*100}}
\fi


% ---- INSERTED FROM DRAWINGS: Level 1 and Level 2 (MCTR II) -------------
Tables~\ref{tab:interfaces-mctr2} and~\ref{tab:interfaces-mctr2-internal} list every interface that crosses the
MCTR II boundary (Level 1 numbering) and every interface between its internal
functional units (Level 2 numbering). The type letter follows the drawing:
M = mechanical, E = electrical, S = signal. Entries marked
\textcolor{red}{TBD} were inferred from the drawing labels and must be
confirmed by the owner of the other party before sign-off.
 
\begin{table}[H]
    \centering
    \caption{External interfaces of MCTR II (F/U \fuMII, Level 1).}
    \label{tab:interfaces-mctr2}
    \footnotesize
    \renewcommand{\arraystretch}{1.25}
    \begin{tabularx}{\textwidth}{@{} l p{2.6cm} p{1.6cm} X @{}}
        \toprule
        \textbf{Interface} & \textbf{From} & \textbf{To} & \textbf{Description and key parameters} \\
        \midrule
        I/F 1.1 M & F/U 1.3 (mounting plate) & F/U \fuMII &
            Bolted mounting of the hoist baseplate (6061-T6, \tbd{mm}) on the
            MCTR I mounting plate. Plate carries the rope reaction of
            \showv{FropeD}{0}{N} (design case). Hole pattern \tbd{}.
            \tbd{confirm against drawing} \\
        I/F 1.3 M & F/U \fuMII.4 (rope) & F/U 1.2 &
            Rope termination to the magnetic load attachment. Rated payload
            \InMassRated\,kg plus attachment \InMassAttach\,kg.
            \tbd{confirm against drawing} \\
        I/F 1.7 E & F/U 1.1 & F/U \fuMII &
            3-phase power to the induction motor from the VSD
            (\InMotorP\,W, \InMotorN\,rpm at \InMotorF\,Hz). Delivered
            frequency \tbd{Hz}. Mains feed to the Brake PSU
            \tbd{not shown in drawing: add interface or confirm it is part
            of I/F 1.7 E}. \\
        I/F 1.8 M & F/U 1.1 (motor) & F/U \fuMII.1 &
            Motor shaft (11\,mm, $4\times16$\,mm keyway) to gearbox input
            shaft (\InShaftD\,mm) through a flexible coupling and a separate
            input stub shaft on its own bearings. Coupling \tbd{part no.}. \\
        I/F 1.9 S & F/U \fuMII.6 & F/U 1.5 &
            HX711 digital output (DOUT/SCK), 24-bit, \InSps\,SPS, logic level
            \tbd{3.3/5\,V}. Load reading used for telemetry and interlock. \\
        I/F 1.10 S & F/U 1.5 & F/U \fuMII.3 &
            Brake release command via relay. Active-high releases; the
            default (reset, watchdog, power loss) is brake engaged, in
            hardware. \\
        \bottomrule
    \end{tabularx}
\end{table}
 
\begin{table}[H]
    \centering
    \caption{Internal interfaces of MCTR II (F/U \fuMII, Level 2).}
    \label{tab:interfaces-mctr2-internal}
    \footnotesize
    \renewcommand{\arraystretch}{1.25}
    \begin{tabularx}{\textwidth}{@{} l p{2.6cm} p{1.6cm} X @{}}
        \toprule
        \textbf{Interface} & \textbf{From} & \textbf{To} & \textbf{Description and key parameters} \\
        \midrule
        I/F \fuMII.1 M & F/U \fuMII.1 & F/U \fuMII.4 &
            Gearbox output shaft to drum hub; M4 cross-bolt through a reamed
            hole. Design torque \showv{TdrD}{1}{N\,m} (\InMassDesign\,kg
            overload with service factor \InServiceF). \\
        I/F \fuMII.2 E & F/U \fuMII.2 & F/U \fuMII.3 &
            48\,V DC supply to the brake coil through a relay, with
            flyback suppression. Coil current \tbd{A}. \\
        I/F \fuMII.3 M & F/U \fuMII.3 & F/U \fuMII.1 &
            Brake on the gearbox input stub shaft (provisional;
            drum shaft is the alternative). Required rating
            \showv{BrakeReqIn}{2}{N\,m} (input shaft) or
            \showv{BrakeReqDrum}{1}{N\,m} (drum shaft), safety factor
            \InBrakeSF. \\
        I/F \fuMII.4 M & F/U \fuMII.4 & F/U \fuMII.5 &
            Rope over a sheave suspended from the load cell. Vertical load
            \showv{Fcell}{0}{N} at rated payload; equal side load on the
            cell from the horizontal rope leg. \\
        I/F \fuMII.5 S & F/U \fuMII.5 & F/U \fuMII.6 &
            4-wire bridge signal, \InCellSens\,mV/V rated output, shielded
            cable, full-scale \showv{SigFS}{1}{mV} against an ADC range of
            \showv{AdcRange}{1}{mV}. \\
        \bottomrule
    \end{tabularx}
\end{table}
 
\noindent\textbf{Open interface questions raised by this table:}
(i) the mains supply to the Brake PSU, required by the system specification to be
distributed by EEII~I, does not appear on the drawings; (ii) the power and
current-sense cable of the magnetic attachment must follow the rope without
carrying tension, which is a cable interface between MCTR II and EEII~II that
is also not drawn; (iii) the drawing numbers the MCTR II sub-units 1.3.x,
which duplicates MCTR I (F/U 1.3), so they are renumbered 1.4.x here.



\subsection{REII: Anton Swanepoel}
\label{subsec:reii-interfaces}
REII Remote Control Dashboard by Anton Swanepoel.

Tables~\ref{tab:reii-external} and~\ref{tab:reii-internal}
define the proposed external and internal interfaces in Figure~3.5.
TBD parameters require agreement.
\cite{nwu_sysrs2026,nwu_reii_scope2026}

\begingroup
\setlength{\tabcolsep}{4pt}
\renewcommand{\arraystretch}{1.15}



\begin{table}[H]
    \centering
    \footnotesize
    \caption{External interfaces of REII (F/U 1.5, Level 1).}
    \label{tab:reii-external}

    \begin{tabularx}{\textwidth}{@{}
        >{\raggedright\arraybackslash}p{2.15cm}
        >{\raggedright\arraybackslash}p{2.7cm}
        >{\raggedright\arraybackslash}p{2.9cm}
        >{\raggedright\arraybackslash}X
    @{}}
        \toprule
        \textbf{Interface} &
        \textbf{From} &
        \textbf{To} &
        \textbf{Description and key parameters} \\
        \midrule

        Operator link\newline (1.5.1) &
        Operator &
        Python dashboard\newline (F/U 1.5.1) &
        Operator commands enter; displayed telemetry and faults return.
        \\[3pt]

        Ground supply\newline (1.5.2) &
        PC supply &
        Ground station and radio\newline (F/U 1.5.2, 1.5.3) &
        PC USB power supplies the UART-connected ground station;
        its radio uses regulated 3.3~V.
        \\[3pt]

        Crane supply\newline (1.5.3) &
        Crane supply &
        Controller and hoist radio\newline (F/U 1.5.5, 1.5.4) &
        Regulated 5~V would feed the S3 board; its 3.3~V
        peripheral rail would supply the radio.
        \\[3pt]

        Inverter control\newline (1.5.4) &
        ESP32-S3 controller\newline (F/U 1.5.5) &
        Drive electronics\newline (F/U 1.1) &
        Six isolated PWM paths plus reserve; current feedback
        below 2.048~V for fault inference.
        \textcolor{red}{[TBD: PWM timing and active levels.]}
        \\[3pt]

        Magnet control\newline (1.5.5) &
        ESP32-S3 controller\newline (F/U 1.5.5) &
        Magnetic attachment\newline (F/U 1.2) &
        Isolated magnet PWM; current feedback below 2.048~V
        for investigating an on-load indicator.
        \textcolor{red}{[TBD: current-signature discrimination,
        calibration and load-support criterion.]}
        \\[3pt]

        Lateral control\newline (1.5.6) &
        ESP32-S3 controller\newline (F/U 1.5.5) &
        Lateral movement\newline (F/U 1.3) &
        PWM and direction commands; quadrature encoder feedback
        for lateral position and speed.
        \textcolor{red}{[TBD: encoder scale and polarity.]}
        \\[3pt]

        Spooling interface\newline (1.5.7) &
        ESP32-S3 controller\newline (F/U 1.5.5) &
        Spooling subsystem\newline (F/U 1.4) &
        REII would power the HX711 (3.3~V/GND), drive SCK
        and read DAT. The proposed optocoupler brake command
        is HIGH to release and LOW to hold; loss of brake
        power would engage the brake. Group links:
        I/F~1.9 load and I/F~1.10 brake.
        \textcolor{red}{[TBD: load calibration and final brake polarity.]}
        \\[3pt]

        Protection status\newline (1.5.8) &
        Local protection &
        ESP32-S3 controller\newline (F/U 1.5.5) &
        The hardware E-stop would remove power from the drive
        electronics independently of software (PHY2).A proposed status signal
        would inform the ESP32-S3 when the E-stop is activated.
        \textcolor{red}{[TBD: E-stop wiring, status signal and
        integration testing.]}
        \\

        \bottomrule
    \end{tabularx}
\end{table}



\begin{table}[H]
    \centering
    \footnotesize
    \caption{Internal interfaces of REII (F/U 1.5, Level 2).}
    \label{tab:reii-internal}

    \begin{tabularx}{\textwidth}{@{}
        >{\raggedright\arraybackslash}p{2.15cm}
        >{\raggedright\arraybackslash}p{2.7cm}
        >{\raggedright\arraybackslash}p{2.9cm}
        >{\raggedright\arraybackslash}X
    @{}}
        \toprule
        \textbf{Interface} &
        \textbf{From} &
        \textbf{To} &
        \textbf{Description and key parameters} \\
        \midrule

        PC serial link\newline (1.5.1.1) &
        Python dashboard\newline (F/U 1.5.1) &
        STM32 ground station\newline (F/U 1.5.2) &
        UART serial requests; replies and diagnostics return
        at 115200~bit/s, 8N1.
        \\[3pt]

        Ground radio SPI\newline (1.5.2.1) &
        STM32 ground station\newline (F/U 1.5.2) &
        Ground radio\newline (F/U 1.5.3) &
        SPI and radio-control signals; received frames and
        radio status return.
        \\[3pt]

        Wireless link\newline (1.5.3.1) &
        Ground radio\newline (F/U 1.5.3) &
        Hoist radio\newline (F/U 1.5.4) &
        Bidirectional nRF24 link carrying commands, results
        and operational telemetry.
        \\[3pt]

        Hoist radio SPI\newline (1.5.4.1) &
        Hoist radio\newline (F/U 1.5.4) &
        ESP32-S3 controller\newline (F/U 1.5.5) &
        Received commands pass to the controller; SPI control
        and reply data return.
        \\

        \bottomrule
    \end{tabularx}
\end{table}

\endgroup

\noindent
Controller logic operates at 3.3~V.
\cite{nwu_s3_board2026}




